\section{Leading forms and compatible further filtrations}
\label{a:filtrations}

The exact minimizer gives one finitely generated graded ring. To improve
that ring by a further minimization, we must preserve its canonical form,
the lc equations, and the values of specified functions on the original
variety. This section proves these compatibility statements. Its final
flattening lemma realizes two successive filtrations by a single real
valuation on the original function field.

\subsection{Full initial fields and leading forms}

All associated graded fields in this section include every homogeneous
degree. Thus, for a valued function field \((K,b)\), we write
\[
                         K_b=\Frac(\gr_b K).
\]
This is not the degree-zero residue field. Its transcendence degree is
\(\operatorname{ratrank}(b)+\trdeg_k\kappa(b)\): choose functions of
independent values and lifts of a residue transcendence basis; after
these choices any further homogeneous element is algebraic, by taking
monomial ratios. In particular \(\trdeg_k K_b=\trdeg_k K\) for an
Abhyankar valuation. In the rationally independent smooth monomial chart
of Lemma~\ref{a:abhyankar}, it is the rational normal-variable field over
the stratum field. Every homogeneous fraction is a stratum coefficient
times a Laurent monomial, and the normal variables are algebraically
independent because their weights are rationally independent. A Gaussian
extension by weighted variables adjoins their independent initial
variables to the previous full initial field.

These descriptions also justify the continuities used below. On a
positive monomial cone, the value of a fixed rational function is locally
a difference of minima on finite supports. Form discrepancy is linear
on the corresponding log smooth chart. Both are therefore continuous.

For normal test configurations with normal central fibre, preservation
of a canonical generator is the top-degree case of
\cite[Definition~3.3 and Proposition~3.5(2)]{XuZhuangForms}.
Here we identify the leading rational top form on the full initial field
and retain prescribed rational-function tests.

\begin{lemma}[Leading differential forms]\label{a:leading-form}
Let \(b\) be a real Abhyankar valuation of a characteristic-zero function
field \(K/k\), and choose \(s_1,\ldots,s_q\in K^*\) whose initials form
a transcendence basis of \(K_b/k\), where \(q=\trdeg_k K\).
If \(\eta=g\bigwedge_{i=1}^q d\log s_i\), then
\begin{equation}
 A_\eta(b)=b(g),\qquad
 \eta_b=\bar g\bigwedge_{i=1}^q d\log\bar s_i.
 \label{eq:A6}
\end{equation}
The second formula is independent of the probes \(s_i\), is equivariant
for actions preserving \(b\), and applies to fractional forms by taking
tensor powers.
\end{lemma}

\begin{proof}
Use the rationally independent normal monomial chart, including the
divisors of the finitely many probes in the monomialization. A basis of
logarithmic normal differentials together with logarithmic differentials
of a separating unit tangential basis has wedge of form discrepancy
zero. Each probe is a unit times a Laurent normal monomial. In this
basis its logarithmic differential has coefficients of nonnegative
value. Their degree-zero parts are exactly the coefficients of its
initial logarithmic differential in the initial variables: a logarithmic
normal derivative of a unit reduces to zero, and a tangential derivative
differentiates its restriction. Lifted etale coordinate derivations give
these assertions in the local ring.

The determinant relating the probe wedge to this fixed wedge has
nonzero reduction, since the initial probes are algebraically
independent and the field has characteristic zero. Its value is
therefore zero. This proves the discrepancy formula. Comparing two
sets of probes by the same determinant computation proves the
independence of the leading form. The construction commutes with field
automorphisms preserving the valuation. Applying it to an integral
tensor power and then dividing discrepancies by that power gives the
fractional-form assertion.
\end{proof}

\subsection{Finite presentations and simultaneous specialization}

We next realize a finitely generated associated graded ring by a weight
family, retaining the rational functions that specify its canonical form
and boundary. Finite generating sets of initials and the corresponding
weight presentations are developed in
\cite[Theorem~2.17, Lemmas~3.2 and~3.4, and Section~3.3]{KavehManon}.
The proof below explains how log Izumi and maximal-ideal powers end the
subduction with an exact polynomial representative in our local setting.

\begin{lemma}[Finite adapted presentations]\label{a:adapted-presentation}
Let \(x\in U\) be an affine klt germ and let \(b\) be a quasi-monomial
valuation centred at \(x\). Suppose \(R_b=\gr_b k[U]\) is finitely
generated. One can choose affine algebra generators \(x_i\), all
vanishing at \(x\), whose initials generate \(R_b\). If
\(J\subset k[X_1,\ldots,X_N]\) is their relation ideal, then
\[
 R_b=k[X_1,\ldots,X_N]/\operatorname{in}_{b(x_i)}J,
\]
where \(\operatorname{in}\) always means the \emph{minimum}-weight
initial ideal. Every regular function has a polynomial representative
attaining its valuation and prescribed initial; the analogous statement
holds for rational functions using fractions.
\end{lemma}

\begin{proof}
Add lifts of finitely many graded generators to any affine generating
set, subtracting constants from the latter. All generator weights are
positive. The minimum of a polynomial relation vanishes on their
initials, giving one inclusion. Conversely, a homogeneous relation
between these initials lifts to a function of strictly higher value.
Subtract a polynomial matching that new initial, and continue with
strictly higher-weight polynomials. Only finitely many monomial weights in these generators
occur below any fixed bound, so the residual value can be made
arbitrarily large. Log Izumi then puts the residual in an arbitrarily
high power of the maximal ideal of the local ring. The same holds in
the affine ring: the power of the closed-point maximal ideal is primary.
That ideal is generated by the \(x_i\), so a polynomial expression in
its sufficiently high power has all monomials of strictly higher weight
than the starting relation. This ends subduction with an actual
polynomial relation of the desired initial. The identical argument,
starting from a function instead of a relation, gives adapted
representatives. Local unit denominators have constant initial, so the
affine and local graded rings agree; taking fractions gives the final
assertion.
\end{proof}

\begin{lemma}[Canonical specialization]\label{a:canonical-specialization}
In Lemma~\ref{a:adapted-presentation}, suppose further that \(R_b\) is
normal and that \(U\) has a canonical generator \(\Omega_U\).
Then \((\Omega_U)_b\) generates the canonical divisor of
\(\Spec R_b\): its divisor is zero. There is a one-parameter weight
family realizing this specialization and preserving initial
identifications on any prescribed finite set of rational functions.
\end{lemma}

\begin{proof}
Choose finite relations whose minima generate the initial ideal in
Lemma~\ref{a:adapted-presentation}. Add adapted representatives of the
finite tests, including full probes and the coefficient \(g\) in
\eqref{eq:A6}. Their ties and strict weight comparisons form a finite
system of rational linear equalities and strict inequalities. Approximate
the generator weights, up to common scaling, by positive integers
\(m_i\) preserving that system. Consider the domain
\[
             \mathcal R=k[s,s^{-m_i}x_i]
                    \subset k[U][s,s^{-1}].
\]
It is flat over \(k[s]\), since it is torsion-free. Its fibre at zero
has the minimum-weight initial presentation for these integer weights.
Indeed substitution and homogenization give one inclusion; conversely,
separating the powers of \(s\) in a homogeneous expression lying in
\((s)\) expresses its upstairs value by terms of strictly higher weight.

The chosen initial relations show that this fibre's ideal contains the
prime presenting \(R_b\). The zero-coordinate section makes the fibre
nonempty. Every minimal prime of the principal ideal \((s)\) has height
one in the finite-type domain \(\mathcal R\), so the fibre has dimension
\(\dim U\). The old initial prime has that same dimension, since its
fraction field is \(K_b\). A strict enlargement would have smaller
dimension. Thus the initial ideals coincide and the special fibre is
exactly \(\Spec R_b\), with the required test initials.

All fibres are normal: the nonzero ones are copies of \(U\), and the
zero fibre is normal by hypothesis. Flatness over the smooth curve
therefore gives normality of the total space. The form
\(\Omega_U\wedge ds\) has zero divisor away from the special fibre.
That fibre is irreducible, reduced and Cartier, so multiplication by a
power of \(s\) gives a rational form with zero divisor on the total
space, hence a canonical generator on its smooth locus.

We identify its relative restriction explicitly. At the generic point
of the special fibre, write the tested probes as
\(s_i=s^{e_i}S_i\) and \(g=s^{e_g}G\), with residues
\(\bar s_i\) and \(\bar g\). Because a factor involving \(ds\)
disappears after wedging with \(ds\),
\[
 \Omega_U\wedge ds
       =s^{e_g}G\bigwedge_i d\log S_i\wedge ds.
\]
The initial logarithmic probe wedge is nonzero. The normalizing power
is consequently exactly \(-e_g\), and the relative restricted form is
\(\bar g\bigwedge_i d\log\bar s_i=(\Omega_U)_b\), up to sign.
At every codimension-one point of the normal special fibre that fibre
is smooth; flatness and geometric smoothness there make the family
smooth over the base. Its relative generator restricts to a generator
there. Thus \((\Omega_U)_b\) has zero divisor on \(\Spec R_b\), as
claimed. This argument uses only codimension-one restriction, and makes
no canonical-sheaf base-change assumption at singular points.
\end{proof}

\begin{lemma}[Simultaneous boundary specialization]
\label{a:boundary-specialization}
In the preceding situation, a rational effective boundary expressed as
a sum of divisors of regular equations specializes to the same sum of
the divisors of their initial equations. This identification agrees with
specializing its prime components, with all multiplicities.
\end{lemma}

\begin{proof}
For a principal ideal, valuation multiplicativity makes its associated
graded ideal principal, generated by the equation's initial. We verify
that component ideals can be preserved at the same time. For each of
the finitely many component ideals, use its polynomial preimage in the
adapted presentation. Its minimum initial ideal maps onto the associated
graded ideal: the adapted representatives in
Lemma~\ref{a:adapted-presentation} prove this also for each element of
the ideal.

Homogenize all these polynomial ideals with one extra variable of
weight zero. Each homogeneous minimum ideal has finitely many generators
that are minima of homogeneous elements of the original homogenized
ideal. Preserve these generators, and the ring's initial ideal and the
equation tests, in one positive integral weight approximation. Inclusion
of the old minima in the new minima is then automatic. There can be no
extra elements, since both homogeneous minimum ideals have the original
Hilbert function, obtained by filtering each finite-dimensional
total-degree piece. Dehomogenization preserves the equality: distinct
terms of one homogeneous polynomial cannot coalesce after setting the
extra variable to one.

The component closures in the resulting weight family have these
specialized ideals, as follows by using the analogous family algebras of
the quotient rings. Rescale each boundary equation by its parameter
power. Its divisor has no vertical component and restricts to the
divisor of the nonzero initial equation. Intersection of codimension-one
cycles with the special-fibre parameter is additive. This can also be
checked at the regular codimension-one points of the normal special
fibre, where the total family is smooth over the base. Decomposing the
horizontal equation divisors into their prime components therefore gives
precisely the same restricted boundary cycles and multiplicities.
\end{proof}

\subsection{Torus fields and flattening on the original field}

We have now identified the specialized canonical form and boundary.
The remaining task is to combine two graded constructions while retaining
that information. Torus invariance separates the first grading from the
second and makes the required combination possible.

The value torus has character lattice generated by the occurring
valuation degrees. Its grading valuation is the real cocharacter that
evaluates these degrees, injectively on this lattice. When other
invariant torus actions descend, we also use their joint effective torus.
For a faithful torus grading on a field generated by homogeneous
functions, choose multiplicative rational variables \(z^\mu\) for a
basis of its character lattice \(M\). Such variables exist because the
occurring characters generate \(M\) as a group. If \(K^T\) is the
invariant field, then
\[
                  K=K^T(z^{\mu_1},\ldots,z^{\mu_r}).
\]
Distinct characters give algebraic independence, and division of every
homogeneous generator by its corresponding torus variable produces an
invariant generator. In particular invariants are generated by ratios
of equal-weight homogeneous elements.

A \(T\)-invariant valuation is Gaussian in these variables. To see the
minimum rule, span the finitely many character terms of a sum by finitely
many torus translates and invert their character matrix. Invariance
shows that each term has value at least that of the sum; the valuation
inequality gives the reverse comparison with the minimum. Formula
\eqref{eq:A6} thus applies using the torus variables and full independent
initials of the restricted valuation. In particular a pure cocharacter
valuation evaluates a homogeneous volume form's discrepancy by its
character.

The next construction refines the successive-degeneration argument of
\cite[proof of Theorem~4.16 and Remark~4.17]{LiXuStable}.
We retain a prescribed finite set of rational-function initials and the
leading top form, and record uniform comparison with the first valuation.

\begin{lemma}[Compatible real flattening]\label{a:flattening}
Let \(b\) be as in Lemma~\ref{a:adapted-presentation}, and let \(T\) be
the value torus of \(R=\gr_b k[U]\), so that \(b\) is injective on its
character lattice. Let \(u\) be a \(T\)-invariant quasi-monomial valuation
centred at the vertex of \(\Spec R\), with \(R'=\gr_u R\) finitely
generated. Any stipulated old grading on \(U\) is assumed to preserve
\(b\) and to act on \(R\) through \(T\).

For arbitrarily small positive real \(\xi\), there is an Abhyankar
valuation \(b_\xi\) on the \emph{actual} field \(k(U)\) with associated
graded ring \(R'\), graded by first weights plus \(\xi\) times second
weights. Its value torus is the joint torus generated by \(T\) and the
value torus of \(u\) on \(R'\). It preserves the stipulated grading.
For any finite set of rational functions, one can preserve both their
double initials and the formula
\[
              b_\xi(F)=b(F)+\xi u(\operatorname{in}_b F).
\]
One can also arrange, on every nonzero regular function,
\[
                    \tfrac12 b(F)\leq b_\xi(F)\leq2b(F).
\]
If \(U,R\) have the canonical generators in
Lemma~\ref{a:canonical-specialization} and both \(b,u\) have their
respective form discrepancies equal to one, the finite tests can be
chosen so that
\[
 A_{\Omega_U}(b_\xi)=1+\xi,\qquad
          (\Omega_U)_{b_\xi}=((\Omega_U)_b)_u.
\]
\end{lemma}

\begin{proof}
First use the lexicographically ordered valuation
\[
           V(F)=(b(F),u(\operatorname{in}_b F)).
\]
On each first-weight slice its second graded calculation is exactly the
one for \(u\). Invariance makes the \(u\)-filtration split over first
degrees. Hence its associated graded on \(k[U]\) is precisely \(R'\),
with the joint grading. Choose lifted joint homogeneous generators,
including the additional generators needed for \(R\) and \(k[U]\).
Take them homogeneous for the stipulated old grading and vanishing at
the point. In the lexicographic subduction procedure there are only
finitely many possible values with bounded first coordinate: the first
graded pieces are finite-dimensional and only finitely many occur up to
that bound. Subduction therefore works as in
Lemma~\ref{a:adapted-presentation}, using log Izumi for the first
coordinate. In particular the lexicographic minimum ideal is the prime
presenting \(R'\), and adapted polynomial and fractional representatives
exist.

Choose \(\xi>0\) small enough to preserve the minima of finitely many
relations generating that initial prime, and avoid all rational ties on
the joint character lattice. The flattened initial ideal contains this
prime. Suppose it contained an extra initial relation. Approximate the
positive flattened generator weights by positive integral weights,
preserving that extra minimum together with the previous finite minima.
The weight-family argument of Lemma~\ref{a:canonical-specialization}
then gives a nonempty special fibre of dimension \(\dim U\), whose
ideal strictly contains the prime of \(R'\), also of dimension
\(\dim U\). This is impossible. Thus the flattened minimum ideal is
exactly the desired prime.

For completeness, the resulting real quotient filtration does define a
valuation on the original ring. Its weight on a function is the maximum
of the minimum weights of all polynomial representatives. Positivity
of the finitely many generator weights makes the weights locally finite.
If the minima of representatives could be arbitrarily high, every
monomial in those representatives would have arbitrarily large ordinary
total degree, putting the function in every maximal-ideal power, hence
making it zero. The maximum is thus
attained for every nonzero function. At an attained maximum, polynomial
minima give the graded quotient by the minimum relation ideal: an
improvement of a representative differs from it by a relation with that
minimum. The quotient being a domain proves multiplicativity, and the
triangle inequality is immediate. Extending by fractions yields
\(b_\xi\). Its full graded field has transcendence degree \(\dim U\),
so it is Abhyankar. Avoidance of rational ties identifies its value
torus with the full joint torus. The homogeneous presentation preserves
the old grading.

To preserve any finite test list, choose its lex-adapted polynomial
representatives and denominators and retain their leading comparisons.
For uniform comparison on \emph{all} regular functions, take \(\xi\)
smaller until every generator's flattened weight lies between half and
twice its first weight. A \(b\)-adapted representative gives
\(b_\xi(F)\geq b(F)/2\). Conversely, a \(b_\xi\)-adapted
representative gives \(b(F)\geq b_\xi(F)/2\). This uses the attained
maximum description separately for the two valuations, and so includes
functions whose initially chosen representatives have cancellation.

Finally choose first-stage probes as follows. Use rational torus
variables and functions of the invariant field whose restricted
\(u\)-initials form a full transcendence basis there. Gaussianity and
Abhyankar equality show that these have full independent first and
double initials: the second-stage torus initials adjoin independent
variables to the full restricted initial field. The invariant probes
are ratios of equal-first-weight homogeneous functions, so they lift
to \(k(U)\), as do the torus variables. Include these probes and the
coefficient of \(\Omega_U\) in their logarithmic basis among the finite
tests. Applying \eqref{eq:A6} first to \(b\), then to \(u\), and then
to the flattening gives coefficient values \(1\), \(1\), and
\(1+\xi\), respectively. The same formula identifies the flattened
leading form with the iterated leading form. This proves the last two
assertions.
\end{proof}

\subsection{A cone proof of the comparison for Cartier-equation boundaries}
\label{a:comparison-cone-proof}

The preceding constructions give a second proof of
Lemma~\ref{a:minimizer-comparison} in the setting actually used later:
\(U\) has a canonical generator, and the klt boundary is a nonnegative
rational combination of regular equations. This proof explains why tests
with nonclosed centres on the stable cone cause no loss of uniformity.
It uses the torus-field description above and the algebraic degeneration
and canonical-form assertions of
Lemmas~\ref{a:canonical-specialization} and
\ref{a:boundary-specialization}.

Let \(F\) be regular affine with \(b(F)>0\); local denominators that are
units have value zero. Degenerate to the stable cone, retaining the
initial equations of the boundary and of \(F\).
The grading \(b\) is still minimizing with the same pair discrepancy,
and the canonical form is the leading form.
Normalize \(A(b)=1\). For an invariant prime-divisor valuation \(E\)
centred on the cone, not necessarily at the vertex, normalize \(A(E)=1\).
On homogeneous elements use \(w_s=b+sE\), and extend by the minimum rule.
Gaussianity in rational torus variables shows this is a quasi-monomial
valuation: its invariant-field restriction is multiplied by \(s\), and
its torus weights are translated by those of \(b\).
It is centred at the vertex since \(w_s\geq b\) on the ring.
The leading-form rule and the homogeneous boundary equations give
\(A(w_s)=1+s\).

Put \(L_*=E(\bar F)/b(\bar F)>0\) and \(s=1/L_*\).
For \(k_m=\lfloor m/(2b(\bar F))\rfloor\), multiply the homogeneous
pieces whose \(b\)-weights lie in
\[
 [\,3m/4-k_m b(\bar F),\ m-k_m b(\bar F)\,)
\]
by \(\bar F^{k_m}\).
They inject into the \(b\)-jets of weight below \(m\).
Their \(b\)-weight is at least \(3m/4\), and their added \(sE\)-weight
is at least \(s k_m E(\bar F)\), so they vanish among the \(w_s\)-jets
for large \(m\). The lost dimension fraction tends to
\(c_p=2^{-p}-4^{-p}>0\). Therefore
\[
 \vol(w_s)\leq(1-c_p)\vol(b),\qquad
 1\leq (1+1/L_*)^p(1-c_p)
\]
by minimality. This bounds \(L_*\) by a constant depending only on \(p\).

Equivariant log resolution shows that invariant prime divisors suffice
to test the homogeneous pair together with \(\bar F\). Undoing the
normalizations, adding \(a\Div(\bar F)\) preserves lc near the vertex
whenever
\[
                      a<\frac{A(b)}{C_p b(F)}
\]
for a dimensional \(C_p\), which may be enlarged from the first proof.
For a strict inequality the cone pair is klt.
Transfer this test through the weight family of
Lemma~\ref{a:canonical-specialization}. Its total space is normal with
Cartier canonical divisor; the special fibre is reduced, normal and
Cartier. All other boundaries are nonnegative rational Cartier-equation
divisors without a vertical component. At codimension one of the
special fibre the family is smooth, and the different is exactly their
Cartier restriction. The ordinary inversion-of-adjunction theorem
\cite[Theorem on p.~1]{Kawakita} applies with coefficient one on the
special fibre. It makes the total pair lc near the vertex.

Away from parameter zero the family and pair are a product. The section
of the original closed point, with all chosen affine generators zero,
specializes to the vertex. Hence the original pair with \(a\Div(F)\)
is lc near that point. Letting \(a\) increase to the displayed bound
gives \eqref{eq:A3}. This argument uses stable degeneration only for
the ordinary klt minimizer, and makes no finite-generation assertion
about an arbitrary lc place.
